% Joint and marginal distributions surface plot used in M3f-Joint-Distributions.tex.
% Precomputed here because the pgfplots surface is slow to render on every lecture build.
% Each overlay below becomes one page of joint-marginal-surface.pdf (pages 1--5 correspond
% to slides 4--8 of the "Joint and marginal distributions" frame).
% Rebuild with:
%   ./build-plots.sh joint-marginal-surface
% (beamer is used so fonts match the slides; the script crops every page to a common box.)
\documentclass[smaller]{beamer}
\usepackage{tikz}
\usepackage{pgfplots}
\pgfplotsset{compat=newest}

\setbeamertemplate{navigation symbols}{}
\setbeamertemplate{footline}{}
\setbeamertemplate{headline}{}

\begin{document}
\begin{frame}[plain]
  \centering
  \pgfplotsset{
    colormap={whitered}{color(0cm)=(white); color(1cm)=(purple!75!blue)}
  }

  \begin{tikzpicture}[scale=.6,
    declare function={mu1=1;},
    declare function={mu2=2;},
    declare function={sigma1=0.5;},
    declare function={sigma2=1;},
    declare function={normal(\m,\s)=1/(2*\s*sqrt(pi))*exp(-(x-\m)^2/(2*\s^2));},
    declare function={bivar(\ma,\sa,\mb,\sb)=
      1/(2*pi*\sa*\sb) * exp(-((x-\ma)^2/\sa^2 + (y-\mb)^2/\sb^2))/2;}]
    \begin{axis}[
      colormap name=whitered,
      width=15cm,
      view={45}{65},
      enlargelimits=false,
      grid=major,
      domain=-1:4,
      y domain=-1:4,
      samples=26,
      xlabel=$x$,
      ylabel=$y$,
      zlabel={$P$},
      colorbar,
      colorbar style={
        at={(1.1,0)},
        anchor=south west,
        height=0.25*\pgfkeysvalueof{/pgfplots/parent axis height},
        title={$P(x,y)$}
      }
      ]
      \addplot3 [domain=-1:4,samples=31, samples y=0, red, thick, smooth] (x,4,{normal(mu1,sigma1)});
      \addplot3 [domain=-1:4,samples=31, samples y=0, blue, thick, smooth] (-1,x,{normal(mu2,sigma2)});
      \addplot3 [surf] {bivar(mu1,sigma1,mu2,sigma2)};

      \only<2->{\draw [black!50] (axis cs:-1,0,0) -- (axis cs:4,0,0);}
      \only<3->{\draw [black!50] (axis cs:0,-1,0) -- (axis cs:0,4,0);}

      \only<5->{\node[blue] at (axis cs:-1,1,0.18) [pin=165:$f(x_2)$] {};}
      \only<4->{\node[red] at (axis cs:1.5,4,0.32) [pin=-15:$f(x_1)$] {};}
    \end{axis}
  \end{tikzpicture}
\end{frame}
\end{document}
